\chapter{Anomaly detection: a model that completes the anomaly}\label{ch:hsldetect}
\def\tbrk{\hfil\penalty50\hfilneg}% local: a line may break before a provenance tag

\begin{onething}
HyperSL rebuilds a pixel's hidden bands from that same pixel's visible bands, so it learns to complete whatever spectrum
it is shown, anomalies included; measured on the 909-object PRISMA bench with the evaluator's own measures, its residual
found 130 of the 909 objects in its top 1\% (random scores find about 93), and even scored with our own residual
Mahalanobis its reconstruction made the score worse than using no reconstruction at all (pAUC 0.083 against 0.128).
\end{onething}

\Cref{ch:recon} built anomaly detection on one rule: \emph{a pixel must never help reconstruct itself}. Indradhanu obeys
it by hiding whole pixels and rebuilding them from their neighbours (\cref{ch:masking}). HyperSL's masking hides bands
of a pixel and rebuilds them from the rest of the same pixel. This chapter works out, on a three-band toy, why that is
the opposite of what a detector needs, and then measures it, because a theory about a nine-million-parameter network is
only a prediction until it is run.

\begin{analogy}
A cloze test, in which a reader fills in the missing words of a sentence. A good reader fills them in for \emph{any}
sentence, however strange, because the rest of the sentence says what is missing: ``the giraffe ordered a \underline{\ \
} of tea'' is completed as easily as ``the man ordered a cup of tea''. If you want to know whether a sentence is odd, a
fluent cloze reader is the wrong judge: it repairs the oddness. Ask instead whether the sentence fits the paragraph it is
in. HyperSL is the fluent cloze reader of spectra; the paragraph is the pixel's neighbourhood.
\end{analogy}

%======================================================================
\section{Two ways to predict a pixel}\label{sec:hsldetect:two}

\begin{definition}[Spatial completion, spectral self-completion]
\emph{Spatial completion}\index{spatial completion} predicts a pixel's spectrum from \emph{other} pixels, its
neighbours. \emph{Spectral self-completion}\index{spectral self-completion} predicts some of a pixel's bands from its
\emph{other bands}. Indradhanu does the first; HyperSL does the second (\cref{tab:hslwhat:trace}, step 5).
\end{definition}

A reconstruction detector scores a pixel by how badly it is predicted (\cref{ch:recon}). For that to flag an anomaly, the
predictor must \emph{fail} on it. Spatial completion fails on an anomaly because its neighbours are background: they
carry no information about the anomaly. Self-completion uses the anomaly's own bands, which carry a great deal.

\subsection{A three-band toy, by hand}

Take spectra of three bands. Ordinary background is \emph{smooth}: the middle band lies halfway between the outer two.
\begin{align*}
\text{background } \vect{b} &= (0.20,\ 0.30,\ 0.40) \tagA,\\
\text{smooth anomaly } \vect{t} &= (0.60,\ 0.50,\ 0.40) \tagA \quad\text{(a bright roof falling with wavelength)},\\
\text{anomaly with a dip } \vect{d} &= (0.60,\ 0.30,\ 0.60) \tagA \quad\text{(an absorption feature in band 2)}.
\end{align*}

\paragraph{The self-completer.} Hide band 2 and predict it from bands 1 and 3 of the same pixel. What predictor would
training on background learn? Fit $\hat x_2=w_1x_1+w_3x_3+c$ by least squares on many background spectra. If every
background spectrum is exactly smooth, $x_2=(x_1+x_3)/2$, then $w_1=w_3=\tfrac12$, $c=0$ gives zero error on every one of
them, and since a squared error cannot be negative, nothing does better. On 2,000 noisy smooth backgrounds the fit gives
$w_1=0.4981$, $w_3=0.5011$, $c=-0.00018$\tagC{} (offsets and slopes drawn at random, noise 0.005\tagA): the average, as
predicted.

\begin{worked}[label={we:hsldetect:self}]{Self-completion of band 2}
With $\hat x_2=(x_1+x_3)/2$:
\begin{align*}
\text{background:}\quad \hat x_2 &= (0.20+0.40)/2 = 0.30, & r_2 &= 0.30-0.30 = 0,\\
\text{smooth anomaly:}\quad \hat x_2 &= (0.60+0.40)/2 = 0.50, & r_2 &= 0.50-0.50 = 0,\\
\text{dip:}\quad \hat x_2 &= (0.60+0.60)/2 = 0.60, & r_2 &= 0.30-0.60 = -0.30 .
\end{align*}
The learned predictor gives $0.4991$ on the smooth anomaly (true 0.50) and $0.5994$ on the dip (true 0.30)\tagC. The smooth
anomaly, three times brighter than the background in band 1, has \emph{the same residual as the background: zero}.
\end{worked}

\paragraph{The spatial completer.} Predict every band of the pixel as the mean of its four neighbours, which are
background.
\begin{worked}[label={we:hsldetect:spatial}]{Spatial completion of the same pixels}
\begin{align*}
\text{background:}\quad \vect{r} &= \vect{b}-\vect{b} = (0,0,0), & \norm{\vect{r}}^2 &= 0,\\
\text{smooth anomaly:}\quad \vect{r} &= (0.60-0.20,\ 0.50-0.30,\ 0.40-0.40) = (0.40,\ 0.20,\ 0),\\
\norm{\vect{r}}^2 &= 0.16+0.04+0 = 0.20,\\
\text{dip:}\quad \vect{r} &= (0.60-0.20,\ 0.30-0.30,\ 0.60-0.40) = (0.40,\ 0,\ 0.20),\\
\norm{\vect{r}}^2 &= 0.16+0+0.04 = 0.20 .
\end{align*}
Both anomalies stand out, by the same amount, whatever their shape\tagC.
\end{worked}

\begin{keypoint}
A self-completer flags only the part of an anomaly it cannot infer from the anomaly itself: a feature narrower than the
spacing of the bands it can see. Everything broad about the anomaly (its brightness, its slope, its overall shape) is
completed, and the residual is the background's.
\end{keypoint}

\subsection{What this means for HyperSL}

HyperSL is a far more flexible self-completer than an average: a transformer that learned, from 200--300 million
spectra, what bands go with what. At mask 0.95 it sees 7 of 149 bands, chosen at random; those 7 carry the anomaly's
overall level and shape, and the other 142 are filled in to match. Three details push it further towards completion.
\begin{itemize}
\item \textbf{Its loss covers visible bands too} (\code{MSE_SAM_loss} on the whole output), so copying what it sees is
rewarded, the identity trap of \cref{sec:recon:identity}.
\item \textbf{The min-max frame pins an anomaly to the edge of $[0,1]$} (\cref{we:hslnorm:anom}), where a sigmoid output is
comfortable.
\item \textbf{It has no neighbours to disagree with.} Even a perfect model of ``what spectra look like'' cannot know that
this spectrum does not belong \emph{here}.
\end{itemize}

\paragraph{The first evidence (2026-09-23).} The matching run of \cref{sec:hslnorm:measured} also measured HyperSL's own
reconstruction error (angle between pixel and reconstruction, mask 0, per-curve scaling) on the ground-truth anomaly pixels
and on random background:

\begin{center}\small
\begin{tabular}{@{}lrrr@{}}
\toprule
scene & anomaly, median angle & background, median angle & AUC (anomaly above background) \\
\midrule
set1 & 3.18° & 6.17° & 0.094 \\
set2 & 4.02° & 5.16° & 0.220 \\
set3 & 3.26° & 4.27° & 0.242 \\
set4 & 2.74° & 3.99° & 0.104 \\
\bottomrule
\end{tabular}
\end{center}
An AUC below 0.5 means the anomalies were reconstructed \emph{better} than the background, on every scene; the mean is
0.165\tagC\tagMr{\texttt{summary.json}, \texttt{recon}}. That run showed HyperSL every band (mask 0), its weakest setting,
and scaled per curve, off its training distribution. The run below fixes both.

%======================================================================
\section{The measurement on the labelled bench (2026-10-31)}\label{sec:hsldetect:run}

For this book the question was put to HyperSL directly, under the conditions it was trained for, and scored exactly as
Indradhanu is scored. The experiment is \file{research/experiments/hypersl-anomaly-bench/} (README, scripts and
\file{results.json}); its score maps are outside the repository in
\file{C:/Users/FITT/Documents/stage1_datasets/hypersl_anomaly_bench/}.

\begin{center}\small
\begin{tabular}{@{}>{\raggedright\arraybackslash}p{24mm}>{\raggedright\arraybackslash}p{112mm}@{}}
\toprule
item & setting \\
\midrule
data & the 4 labelled PRISMA L2D scenes, 909 objects; validity and object maps exported from the evaluator's own
\code{PrismaBench} (\file{step0_export_masks.py}), so every method is counted on the same pixels \\
HyperSL input & the 149-band PRISMA axis of every earlier HyperSL experiment; each band min-max scaled over the scene's
valid pixels, as in \file{mutihsi.py} \\
model & \code{10_base_mask95} (sha256 \code{3d5f1914...a74a68c}), mask ratio 0.95 (7 of 149 bands visible), 4 random
masks per pixel (seeds 1000--1003), fp16 on an RTX 4000 SFF Ada in WSL (\file{step1_hypersl_scores.py}) \\
reconstruction & each band's prediction averaged over the draws in which it was \emph{hidden}, so every residual is a
prediction of a band the encoder did not see \\
HyperSL scores & \textbf{L2}: mean squared residual in min-max units; \textbf{loss}: its own training objective
$\text{MSE}+(1-\cos)$ per pixel; \textbf{RM}: the product's residual Mahalanobis (30 components, ring 21/61) on the
residual converted back to reflectance \\
baselines & \textbf{FTH3 epoch 5}, Indradhanu's release candidate (its cached score maps from the 2026-10-29 evaluation,
found by the evaluator's cache key); \textbf{PCA-RX}, the evaluator's \code{pca_rx} (residual Mahalanobis with a
\emph{zero} reconstruction); \textbf{MNF-RX}, the product's \code{MNFCompressionDetector}, 10 components \\
cloud & every method also scored with FTH3 epoch 5's thick-cloud pixels set to 0, the release candidate's rule \\
measures & \code{app.evaluation.metrics} and \code{significance} (\cref{ch:measures,ch:bootstrap}), via
\file{step2_evaluate.py} \\
\bottomrule
\end{tabular}
\end{center}

\paragraph{Three checks before the results.} FTH3's two rows reproduce its evaluation report exactly (pAUC 0.1749 and
0.1892, 422 and 449 objects)\tbrk\tagM{research/experiments/indradhanu-finetune/FTH3-heavy-joint-more2/eval_report.md}.
PCA-RX gives 0.1281, the stage-3 run's 0.128\tbrk\tagM{docs/lld/evaluation.md}. And the masking did what it should:

\begin{worked}[label={we:hsldetect:never}]{How many bands should never be hidden in 4 draws?}
In one draw a band is visible with probability $7/149$ (7 of 149 tokens kept). In 4 independent draws it is visible
every time with probability
\begin{align*}
\left(\frac{7}{149}\right)^4 &= 0.046980^4 = 0.0022071^2 = 4.8713\times10^{-6}.
\end{align*}
Set1 has 992,708 valid pixels of 149 bands:
\begin{align*}
992{,}708\times149 &= 147{,}913{,}492 \text{ (pixel, band) pairs},\\
147{,}913{,}492\times4.8713\times10^{-6} &= 720.5 .
\end{align*}
The run logged 706\tagM{research/experiments/hypersl-anomaly-bench/results.json}: the expected count, within chance
($\sqrt{720}\approx27$)\tagC. For those 706 entries the script used the last draw's reconstruction, made with the band
visible.
\end{worked}

\subsection{The results}

\begin{table}[htbp]
\centering\small
\begin{tabular}{@{}lrrrrr@{}}
\toprule
method & caught top 1\% & caught top 0.1\% & pAUC & ROC-AUC & TPR at 0.001 \\
\midrule
\emph{chance} & 10.2\% & & & & \\
FTH3 ep.\ 5, RM & 422 (46.4\%) & 13.9\% & 0.1749 & 0.866 & 0.058 \\
PCA-RX (no model) & 361 (39.7\%) & 8.1\% & 0.1281 & 0.839 & 0.029 \\
\textbf{HyperSL RM} & 350 (38.5\%) & 5.4\% & 0.0832 & 0.757 & 0.020 \\
MNF-RX & 249 (27.4\%) & 3.4\% & 0.0750 & 0.749 & 0.019 \\
\textbf{HyperSL L2} & 130 (14.3\%) & 4.0\% & 0.0175 & 0.624 & 0.007 \\
\textbf{HyperSL loss} & 63 (6.9\%) & 0.9\% & 0.0039 & 0.501 & 0.000 \\
\midrule
\multicolumn{6}{@{}l}{\emph{cloud excluded by FTH3 epoch 5's thick-cloud mask}} \\
FTH3 ep.\ 5, RM & 449 (49.4\%) & 16.0\% & 0.1892 & 0.869 & 0.065 \\
PCA-RX & 377 (41.5\%) & 8.5\% & 0.1390 & 0.842 & 0.031 \\
\textbf{HyperSL RM} & 369 (40.6\%) & 6.7\% & 0.0947 & 0.761 & 0.024 \\
MNF-RX & 280 (30.8\%) & 5.2\% & 0.0966 & 0.753 & 0.024 \\
\textbf{HyperSL L2} & 137 (15.1\%) & 4.3\% & 0.0182 & 0.628 & 0.007 \\
\textbf{HyperSL loss} & 63 (6.9\%) & 0.9\% & 0.0039 & 0.503 & 0.000 \\
\bottomrule
\end{tabular}
\caption{HyperSL on the 909-object PRISMA bench, beside Indradhanu FTH3 epoch 5 and two classical detectors. pAUC is up to
FPR 0.01; pixel measures are means over the four scenes; caught is the share of all 909 objects with a pixel within 1\,px
of the top 1\% or 0.1\%\tagM{research/experiments/hypersl-anomaly-bench/results.json}.}
\label{tab:hsldetect:main}
\end{table}

\begin{figure}[htbp]
\centering
\begin{tikzpicture}
\begin{axis}[xbar, bar width=6pt, width=0.62\linewidth, height=62mm, xmin=0, xmax=0.23, xlabel={pAUC (FPR $\le$ 0.01)},
  xtick={0,0.05,0.1,0.15,0.2}, xticklabel style={/pgf/number format/fixed,/pgf/number format/precision=2},
  scaled x ticks=false,
  ytick={0,1,2,3,4,5}, yticklabels={HyperSL loss,HyperSL L2,MNF-RX,HyperSL RM,PCA-RX (no model),FTH3 ep.\ 5},
  y tick label style={font=\small}, enlarge y limits=0.12, legend style={at={(0.98,0.03)},anchor=south east,font=\scriptsize},
  nodes near coords, nodes near coords style={font=\tiny}, every node near coord/.append style={/pgf/number format/fixed,/pgf/number format/precision=3}]
\addplot[fill=maskgrey] table[y=pos, x=pauc] {figdata/hsl_detect_bar.dat};
\addplot[fill=spec!70] table[y=pos, x=paucx] {figdata/hsl_detect_bar.dat};
\legend{plain, cloud excluded}
\end{axis}
\end{tikzpicture}
\caption{pAUC on the PRISMA bench. HyperSL's own residual (L2, loss) sits near zero; scored with residual Mahalanobis it
lands below PCA-RX, which uses no model at all. Data \texttt{figdata/\allowbreak hsl\_detect\_bar.dat}.}
\label{fig:hsldetect:bar}
\end{figure}

\paragraph{1. HyperSL's own residual is close to chance.} Its training loss as a score catches 63 objects, \emph{fewer}
than the 10.2\% (about 93) that random scores would catch\tagC, with a ROC-AUC of 0.501. Its mean squared residual does
somewhat better, 130 objects and pAUC 0.0175, a tenth of FTH3's. This is the toy of \cref{we:hsldetect:self} at full
scale: the anomalies are completed.

\paragraph{2. Our scorer rescues part of it, but not past ``no model''.} Feeding HyperSL's residual to the product's
residual Mahalanobis gives 350 objects and pAUC 0.0832. Most of that is the scorer's work, as the next comparison shows.

\paragraph{3. The decisive comparison: HyperSL's reconstruction against none.} PCA-RX \emph{is} residual Mahalanobis on
the cube with a zero reconstruction. So ``HyperSL RM minus PCA-RX'' isolates what subtracting HyperSL's reconstruction
does, everything else equal. Paired object bootstrap, 2,000 draws:

\begin{center}\small
\begin{tabular}{@{}lll@{}}
\toprule
difference (B $-$ A) & pAUC [95\% interval] & caught top 1\% [95\% interval] \\
\midrule
HyperSL RM $-$ PCA-RX & $-0.0450$ [$-0.0589$, $-0.0318$] & $-1.2$ [$-4.0$, $+1.5$] points \\
HyperSL RM $-$ PCA-RX, cloud excl. & $-0.0443$ [$-0.0577$, $-0.0320$] & $-0.9$ [$-3.5$, $+1.9$] points \\
PCA-RX $-$ FTH3 ep.\ 5 & $-0.0467$ [$-0.0604$, $-0.0327$] & $-6.7$ [$-8.8$, $-4.5$] points \\
HyperSL RM $-$ FTH3 ep.\ 5 & $-0.0917$ [$-0.1060$, $-0.0776$] & $-7.9$ [$-10.8$, $-5.2$] points \\
HyperSL RM $-$ FTH3 ep.\ 5, cloud excl. & $-0.0944$ [$-0.1094$, $-0.0800$] & $-8.8$ [$-11.7$, $-6.2$] points \\
MNF-RX $-$ FTH3 ep.\ 5 & $-0.0999$ [$-0.1170$, $-0.0826$] & $-19.0$ [$-22.0$, $-16.3$] points \\
\bottomrule
\end{tabular}
\end{center}
\textbf{Subtracting HyperSL's reconstruction lowers pAUC by 0.045, an interval well clear of zero}, and leaves the object
count statistically unchanged. Indradhanu's reconstruction does the opposite: against the same PCA-RX it \emph{adds} 0.047
(the third row, read with its sign reversed). The two models' residuals move the same scorer by almost the same amount in
opposite directions\tagC. That is the measured form of \cref{sec:hsldetect:two}: a spatial completer leaves the anomaly in
the residual; a self-completer takes it out.

\begin{worked}[label={we:hsldetect:sym}]{Reading the bootstrap rows}
FTH3 against PCA-RX is the third row reversed: $+0.0467$ with interval $[+0.0327,+0.0604]$. HyperSL against PCA-RX is
$-0.0450$ with interval $[-0.0589,-0.0318]$. Add the two point values:
\begin{align*}
0.0467+(-0.0450) &= 0.0017 .
\end{align*}
Within $0.002$ of zero: on this bench, what Indradhanu's residual adds to the cube's own statistics is about what HyperSL's
residual removes\tagC. Both intervals exclude zero, so both effects are resolved.
\end{worked}

\paragraph{Per scene and by contrast tier.}

\begin{center}\small
\begin{tabular}{@{}lrrrrrrr@{}}
\toprule
& \multicolumn{4}{c}{pAUC per scene} & \multicolumn{3}{c}{caught top 1\% by tier} \\
\cmidrule(lr){2-5}\cmidrule(lr){6-8}
method & set1 & set2 & set3 & set4 & easy & medium & subtle \\
\midrule
FTH3 ep.\ 5, RM & 0.2616 & 0.1874 & 0.0543 & 0.1964 & 59.2\% & 41.0\% & 39.0\% \\
PCA-RX & 0.1734 & 0.1171 & 0.0822 & 0.1400 & 48.7\% & 35.4\% & 35.0\% \\
HyperSL RM & 0.0875 & 0.0966 & 0.0286 & 0.1199 & 52.0\% & 33.8\% & 29.7\% \\
MNF-RX & 0.0871 & 0.1267 & 0.0196 & 0.0664 & 34.9\% & 25.2\% & 22.0\% \\
HyperSL L2 & 0.0415 & 0.0225 & 0.0040 & 0.0021 & 21.1\% & 10.8\% & 11.0\% \\
HyperSL loss & 0.0021 & 0.0117 & 0.0019 & 0.0000 & 7.9\% & 4.9\% & 8.0\% \\
\bottomrule
\end{tabular}
\end{center}
HyperSL RM is below PCA-RX and below FTH3 on every scene. It holds up best on the easy (high-contrast) tier, 52.0\% against PCA-RX's 48.7\%, and falls furthest on the
subtle one, 29.7\% against 35.0\%: the low-contrast targets are the ones a self-completer rebuilds most
completely\tagM{research/experiments/hypersl-anomaly-bench/results.json}.

\paragraph{Cost.} 428--462\,s of GPU time per scene for about a million pixels and four draws, 1,773\,s for the bench: 438
microseconds per pixel on a 20\,GB laptop GPU\tagC. For comparison, Indradhanu Heavy takes 454\,s per PRISMA scene on 16
CPU threads in the product worker\tbrk\tagMr{commit \texttt{72a767c} message}.

\begin{sourcenote}[: what this run did not try]
HyperSL read the 149-band axis of the earlier experiments, not the 186-band native axis the evaluator gives Indradhanu,
and it kept no water bands, unlike its pretraining. Mask ratio and draw count were not tuned (0.95, the training ratio,
and 4). These choices follow the model's training and the team's earlier runs; none was tuned on the bench, so none can
have flattered the comparison either way. A different band axis could move HyperSL's numbers a little; it cannot remove
self-completion, which is in the masking.
\end{sourcenote}

%======================================================================
\section{HyperSL as a false-alarm filter}\label{sec:hsldetect:filter}

A gentler use was tried on 2026-09-26: keep Indradhanu as the detector, and use HyperSL's embeddings of each candidate
pixel and of its ring to tell real targets from false alarms (\file{research/experiments/indradhanu-fp-filter/}, scripts in
commit \code{74df314}; results outside the repository in
\file{C:/Users/FITT/Documents/stage1_datasets/indradhanu_fp_filter/step3/results.json}). Candidates are the top 2\% of
Indradhanu's residual-Mahalanobis score; a small classifier is trained on three scenes and tested on the fourth; the
headline is the share of false-alarm pixels removed while keeping 90\% of real objects. ``Oracle'' sets that cut on the
test scene's own labels (an upper bound); ``honest'' sets it on the training scenes, as deployment would.

\begin{center}\small
\begin{tabular}{@{}lrrrrr@{}}
\toprule
features (one hidden layer) & set1 & set2 & set3 & set4 & mean \\
\midrule
\multicolumn{6}{@{}l}{\emph{oracle: share of false-alarm pixels removed at 90\% object recall}} \\
none (a stricter Indradhanu cut) & 0.339 & 0.128 & 0.198 & 0.262 & 0.232 \\
detector scores only & 0.059 & 0.219 & 0.223 & 0.249 & 0.188 \\
HyperSL pixel + ring embeddings & 0.615 & 0.464 & 0.188 & 0.318 & 0.396 \\
\midrule
\multicolumn{6}{@{}l}{\emph{honest: object recall reached / false-alarm pixels removed}} \\
none & 1.00 / 0.00 & 0.62 / 0.56 & 1.00 / 0.00 & 0.86 / 0.31 & \\
HyperSL embeddings & 0.72 / 0.77 & 0.42 / 0.87 & 0.45 / 0.50 & 0.69 / 0.59 & \\
\bottomrule
\end{tabular}
\end{center}
\tagMr{\texttt{indradhanu\_fp\_filter/step3/results.json}, configurations \texttt{baseline}, \texttt{scores/B},
\texttt{hypersl/B}; means \tagC}

This is HyperSL's best showing in the whole part, and it deserves to be stated plainly: with the cut chosen in hindsight,
its embeddings removed 40\% of false-alarm pixels at 90\% object recall against 23\% for a stricter cut, so they carry
information about what a target looks like next to its ring. But the honest cut, the only one a deployment has, did not
transfer between scenes: it kept between 42\% and 72\% of the objects instead of 90\%. The filter was not adopted.

\begin{keypoint}
A model that rebuilds a pixel from that pixel cannot tell you the pixel is surprising. Scored as a detector, HyperSL's
residual is near chance; scored with our own residual Mahalanobis, its reconstruction makes the score worse than using no
reconstruction at all.
\end{keypoint}

\begin{drift}
The bench run is of 2026-10-31 (\file{research/experiments/hypersl-anomaly-bench/}); FTH3's maps are from the 2026-10-29
evaluation; the false-alarm-filter numbers are from 2026-09-26 and were computed on candidates from the Indradhanu of that
date. A new HyperSL release, a different band axis, or a new release candidate would move the rows.
\end{drift}

%======================================================================
\begin{confusions}
\confusion{HyperSL hides 95\% of the bands, so it cannot copy.}{It cannot copy those bands, but the 7 it sees carry the
anomaly's level and shape, and the other 142 are completed to match. Hiding bands is not hiding the pixel.}
\confusion{A residual Mahalanobis of 350 objects shows HyperSL detects anomalies.}{The same scorer on the raw cube, with
no model, catches 361 and has a higher pAUC. The objects are the scorer's, not HyperSL's.}
\confusion{HyperSL's loss catching 63 objects is a little worse than chance, so it is useless but harmless.}{Below chance
means it ranks anomalies \emph{lower} than background: it fits them better. That is self-completion, not noise.}
\confusion{Indradhanu would have the same problem with spectral masking.}{It would, which is why spectral self-completion
is listed as rejected in its design (\file{docs/lld/spectral-stem.md}) and band-list augmentation is allowed only because a
hidden pixel's own bands are never an input (\cref{sec:specdecoder:traps}).}
\confusion{The false-alarm filter result shows HyperSL helps.}{Only with a cut chosen on the test scene's labels. With an
honest cut it lost up to 58\% of the objects.}
\end{confusions}

\begin{checkyourself}
\item\label{ex:hsldetect:1} A third anomaly is $(0.50,\ 0.45,\ 0.40)$. Compute its self-completion residual in band 2 and its
  spatial residual and squared length against $\vect{b}=(0.20,0.30,0.40)$.
\item\label{ex:hsldetect:2} Set2 has 1,050,895 valid pixels. How many (pixel, band) pairs do you expect never to be hidden in 4
  draws? The run logged 736.
\item\label{ex:hsldetect:3} An object covers 5 pixels after the 1-pixel dilation. What is its chance of being caught in the top
  1\% by random scores? (\Cref{ch:measures}: $1-(1-q)^m$.)
\item\label{ex:hsldetect:4} Does the bootstrap show HyperSL RM catching fewer objects than PCA-RX? Does it show a lower pAUC?
\item\label{ex:hsldetect:5} Why can a score catch \emph{fewer} objects than chance? What does that say about how HyperSL ranks
  anomalies?
\item\label{ex:hsldetect:6} At 438 microseconds per pixel with 4 draws, how long would a 1,000 by 1,000 scene take with 4 draws,
  and roughly with 1 draw?
\end{checkyourself}

\nextchapter{ch:hslunmix}
